\section{NN-S001: fixed fibres, rather than an automatic Boolean fixed part}
\label{sec:NN-S001}

\paragraph{Inputs.} This is a general form of corpus 01, Section 3
(semimodules as diagrams of modules), corpus 02's corresponding repair,
and corpus 16, Theorem \texttt{thm:fixed-locus-unit}. No E, A, or M is used.

\begin{theorem}[NN-S001]
Let $R$ be a nonzero commutative unital ring, $S=G(R)$, and $M$ an
$S$-semimodule. Write $e=0_R$, $B=eM$, and $M_b=\{m:em=b\}$.
Then $B$ is a join semilattice, $M$ is the disjoint union of the
$R$-modules $M_b$, whose respective zero elements are $b$, and
\[
 \operatorname{Fix}_M(u)=\coprod_{b\in B}\ker(u-1:M_b\longrightarrow M_b)
 \qquad(u\in R^\times).
\]
In particular $\operatorname{Fix}_M(-1)=\coprod_b M_b[2]$.
This is exactly the skeleton $B$ if and only if every fibre has no
nonzero $2$-torsion. The assertion need not hold for an arbitrary
$G(R)$-semimodule even when $R=\mathbb Z$.
\end{theorem}

\begin{proof}
The scalar identities $e+e=e$, $1+e=1$, and $1+(-1)=e$ imply
$b+b=b$, $m+em=m$, and $m+(-1)m=em$. Thus $B$ is an idempotent
commutative monoid, with join its addition. In $M_b$, addition is closed,
$b$ is an identity, and $(-1)m$ is the inverse of $m$. Multiplication
by a supported scalar $r\in R$ preserves $M_b$ because $er=e$.
The semimodule axioms now give precisely the $R$-module axioms in
each fibre, including the action of $0_R$ as its zero element $b$.
The displayed formula is the equation $um=m$ inside this module.
For a counterexample let $M=\mathbb Z/2\mathbb Z$ acted on through
$G(\mathbb Z)\to\mathbb Z$. It has the single skeleton element $0$,
but both module elements are sign-fixed. In contrast every fibre of
$G(R)^d$ is a coordinate module $R^I$; its sign-fixed part is its
Boolean skeleton when $2$ is a non-zero-divisor in $R$.
\end{proof}

\paragraph{Consequence and remaining gate.}
Replacing $\mathbb Q_\ell$ by arbitrary $G(R)$-semimodules in the
sign-only argument does not automatically make the entire fixed part
Boolean. One must specify torsion-free linear fibres and then verify
the trace and effective-spectrum hypotheses on each chosen realization.
This corrects the unrestricted coefficient assertion in prompt 10(a);
it yields no parity or satisfiability algorithm.

\paragraph{Geometric dictionary for NN-S001.}
The skeleton indexes a diagram of ordinary modules, hence of
quasi-coherent sheaves when the modules are localized over $\operatorname{Spec}R$.
The fixed fibre is the kernel of the endomorphism $u-1$ of that sheaf.
This is an exact module calculation, not a claim that the skeleton is
an additional vector-space summand. The ordinary affine-module
dictionary is EGA I, Section 1.3; exact edition locators are recorded
in the source register. The nonadditive skeleton lies outside that
ordinary quasi-coherent category.

\section{NN-S002: the sign-fixed scheme has a characteristic-two fibre}
\label{sec:NN-S002}

\begin{theorem}[NN-S002]
The fixed-point scheme of sign on $\mathbb A^d_{\mathbb Z}$ is
\[
 F_d=\operatorname{Spec}\mathbb Z[x_1,\ldots,x_d]/(2x_1,\ldots,2x_d).
\]
For $d>0$ it is not flat over $\mathbb Z$. Its characteristic-zero
fibre is the origin, whereas its characteristic-two fibre is the
whole affine space. Consequently the reduction of
$\operatorname{Fix}_{G(\mathbb Z)}(-1)=\{\tau,e\}$ is a proper subset
of $\operatorname{Fix}_{G(\mathbb F_2)}(-1)=G(\mathbb F_2)$.
\end{theorem}

\begin{proof}
For every commutative ring $T$, sign fixes a vector $(a_i)$ exactly
when $2a_i=0$ for all $i$. The displayed coordinate algebra represents
this functor. Tensoring it with $\mathbb Q$ or $\mathbb F_2$ gives
the two claimed fibres. The nonzero element $x_1$ in the integral
coordinate algebra is killed by $2$; a flat $\mathbb Z$-module has
no such torsion, since tensoring $0\to\mathbb Z\xrightarrow{2}\mathbb Z$
must preserve injectivity. Over $\mathbb Z$ the scalar fixed equation
has only amplitude zero, while over $\mathbb F_2$ sign is identity.
The external element $\tau$ is fixed in both cases, proving the
last assertion.
\end{proof}

\paragraph{Geometric dictionary and scope.}
This theorem is already a scheme-theoretic statement: it is the
equalizer of identity and the order-two automorphism of the vector
group. The equalizer itself \emph{does} commute with base change,
as its displayed equations show. What fails is replacing it by the
zero section using its characteristic-zero points and then reducing.
The coordinate calculation provides the full proof. The rank-one
example is already recorded in SGA 3, Expos\'e I, Remark 4.7.1.2,
printed page 29 of the corrected edition dated 13 October 2024
(the editorial addition is identified by its note 39). The new
application here is to the split-semimodule claim; the classical
fixed-scheme phenomenon is not claimed as new. Equalizers and base
change use SGA 3, I, Sections 1.2 and 1.6.
For an $R$-module the analogous issue is the kernel of multiplication
by $2$; its change after tensoring is measured by the associated
$\operatorname{Tor}_1$ sequence. No assertion about a finite-field
Frobenius trace follows merely from the sign-fixed Boolean pair.

\section{NN-S003: additive linearization erases the support skeleton}
\label{sec:NN-S003}

\begin{theorem}[NN-S003]
Every additive map from an idempotent commutative monoid to an
abelian group is zero. Thus the universal ring reflection of $G(R)$
is $R$, and extension of scalars along $G(R)\to R$ kills every
semimodule on which $e$ acts as identity. In particular
\[
 R\otimes_{G(R)}\mathbb B=0.
\]
Even on free rank-one split modules, Boolean Frobenius data do not
determine a Weil weight: for each integer $j$, multiplication by
$q^j$ on $G(\mathbb Q_\ell)$ has the same Boolean action, but its
amplitude realization has eigenvalue $q^j$ and weight $2j$.
\end{theorem}

\begin{proof}
If $x+x=x$, then $f(x)+f(x)=f(x)$ in the group, so $f(x)=0$.
Any map from $G(R)$ to a ring therefore identifies $e$ with the
additive identity, which also receives $\tau$. The remaining
supported scalars have exactly the ring operations of $R$, proving
the universal property. In the tensor product, for $em=m$ one has
$r\otimes m=r\otimes em=re\otimes m=0$, since $e$ maps to zero
in $R$. Finally $q^j$ is a nonzero scalar, so its Boolean support
action is identity for every $j$. On the ordinary one-dimensional
realization every complex embedding of its eigenvalue has absolute
value $q^j$. This gives distinct weights with identical skeletons.
\end{proof}

\paragraph{Purity consequence.}
If a finite Boolean set has a Frobenius permutation, its separate
free linearization on the atoms is a permutation representation;
its eigenvalues are roots of unity and it has weight zero. This
linearization is not additive for union: the vector attached to
$A\cup A$ equals that attached to $A$, not twice that vector.
It therefore does not contradict the theorem. Nor is scalar
group completion the same construction as the split $K_0$ of a
category of finite Boolean objects, whose direct-sum ranks can survive.

The correct conditional transfer is fibrewise: supply a finite
effective Frobenius spectrum, the path-count identities and bounds
of GCT\_NSRH F03 (or F04/F17), at every required embedding, and the
dual spectrum. That cited criterion then gives its stated weight
bound and, with duality, purity. The existence of $G(R)$, its fixed
pair, and A provide none of those missing realizations. The same
support is compatible with every Tate weight, so weight zero on
the skeleton cannot force PH1 or positivity of a canonical formula.

\paragraph{Geometric dictionary for NN-S003.}
The amplitude rank-one examples are Tate Weil objects; their
Frobenius eigenvalues give the standard weights of Weil II,
Definition 1.2.2. The distinction between additive group completion
and the Grothendieck group of an exact category is essential in
SGA 6, Expos\'e IV (the definition of $K_0$ and its functorial
operations). No realization functor to that exact category is
constructed by a semiring support map.

\section{NN-S004: arbitrarily large support geometry with the same zeta}
\label{sec:NN-S004}

\paragraph{Background.} The construction and ideal classification
of $G_L(R)$ are corpus 16, Theorems \texttt{thm:lattice-normal-form},
\texttt{thm:lattice-ideal-classification}, and
\texttt{thm:lattice-prime-classification}. They are inputs to the
following zeta consequence, not new constructions in this run.

\begin{theorem}[NN-S004]
Let $L$ be a nontrivial bounded distributive lattice and $R$ an
infinite commutative unital ring. In $G_L(R)$ give an ideal its Bourne quotient norm.
Every finite-norm ideal is an amplitude ideal
$I_J=Z_L\cup\{(r,1_L):r\in J\}$, and its quotient is $R/J$.
Conversely each finite quotient $R/J$ gives this ideal and norm.
Moreover $I_JI_K=I_{JK}$. In particular for every number field $K$
and every such $L$,
\[
 \zeta^{\rm cl}_{G_L(\mathcal O_K)}(s)=\zeta_K(s).
\]
For finite $L$ of cardinality $\ell$, the separated quotient has
cardinality $N(J)+\ell-1$ and the separated sheet is
\[
 \sum_{0\ne J\subseteq\mathcal O_K}(N(J)+\ell-1)^{-s}.
\]
Thus the classical sheet cannot reconstruct the cardinality, height,
or branching of the support lattice, even when its RH is supplied by A.
\end{theorem}

\begin{proof}
A pure-support ideal lies in the kernel of the amplitude map to
$R$. That map therefore descends to its Bourne quotient and is
still surjective. Its quotient norm is infinite. For $I_J$, the
amplitude map modulo $J$ gives a surjection onto $R/J$.
All zero-fibre elements become equivalent to the external zero:
they lie in $I_J$, hence are Bourne-equivalent to zero by adding
the element itself on the other side. Supported scalars differing
by $J$ are equivalent by the ordinary additive equations, and the
map to $R/J$ shows that no other pairs are equivalent. This proves
the quotient claim and exhausts finite-norm ideals by the cited
classification. Products of amplitude ideals have amplitude $JK$;
they contain every $z_\lambda$ because $e_Lz_\lambda=z_\lambda$
with both factors in the two ideals. Conversely every product
amplitude is in $JK$, proving equality. For $R=\mathcal O_K$ the
finite-index ring ideals are exactly the nonzero ideals, with
the usual multiplicative norms. The classical Dirichlet series
is therefore identical term by term to $\zeta_K$.

In a separated quotient $G_L(R/J)$ the zero-amplitude fibre has
$\ell$ elements and each of the other $|R/J|-1$ amplitudes has
one, proving its cardinality. Finally choose chains of arbitrarily
large finite length, or Boolean lattices of arbitrarily large rank.
These have different heights or branching but give the identical
classical Dirichlet series. For the chain with $n+1$ elements,
corpus 16, Corollary \texttt{cor:lattice-split-dimension}, gives
$\dim G_L(\mathbb Z)=n+1$, which is unbounded.
\end{proof}

\paragraph{Geometric dictionary and complexity impact.}
The arithmetic scheme seen by these finite Bourne quotients is
exactly $\operatorname{Spec}\mathcal O_K$; ring reflection kills
the full idempotent lattice. The extra lattice-prime stratum in
corpus 16's spectral ordinal sum is a semiring spectrum, not an
extra stratum of the ordinary EGA scheme. For abelian $K$, A and
abelian Artin factorization give RH for the common classical
sheet. This equality of analytic invariants imposes no reconstruction
of the arbitrarily chosen support lattice. This is an explicit
failure of a proposed inference from the invariant, not a logical
independence proof about P versus NP. It supplies no efficient
algorithm for an implicitly presented support lattice.
